\section{Source-labelled prism regions attached to cut cores}
\label{sec:prism-interfaces}

Whole prismatic cut pieces do not describe the parallelity that remains
inside a marked core. A primitive connected cut can have no whole prismatic
components and still contain many locally parallel discs. This continuation
retains their face-connected prism regions and the exact source faces and
ranks where they meet exceptional chambers. The output is a geometric
incidence representation and a verified partition of interface ports.
It is not yet a full parallelity-bundle or core-handle reconstruction.

\subsection{Remove exceptional points without losing face locations}

Start with the owned chamber model of Section~\ref{sec:cut-complement}.
Delete its at most $6T$ exceptional chamber points. On the remaining
points use the order-preserving compression
\[
             \pi(x)=x-\#\{e\text{ exceptional}:e<x\}.
\]
For each source face region, retain its original source/target tetrahedra,
faces, corner and rank interval. Its two chamber maps are affine with slopes
$+1$ or $-1$. Split the rank interval at every preimage of an exceptional
point and immediately after it. Each resulting part has one of three roles:
\begin{enumerate}
\item Both chambers are prisms: retain the compressed interval pairing and
  one face-band contribution on its source endpoint family.
\item Exactly one chamber is exceptional: retain a singleton interface port
  with its source face, corner and rank, core-cell identity, side and compressed
  prism point.
\item Both chambers are exceptional: retain their source-labelled core-face
  contact. No prism identification is added.
\end{enumerate}

Because no remaining piece crosses a deleted point, compression preserves
its interval isometry. The domain/range normalization used by the orbit
kernel does not erase source-side information: port positions are computed
from the original affine face maps before normalization. Self-face gluings
are included. A self pairing that changes no connectivity still represents
a face band and must contribute to the ribbon profile.

\begin{proposition}[Bounded source ports and exact face-region partition]
The construction has at most $6T$ exceptional cells and $24T$ core/prism
face ports. Its interval orbits are exactly the components of the prism
face-adjacency graph. Every component touching a core-facing port is
identified by the set of those port indices. Components touching no such
port are exactly the whole prismatic cut pieces of the preceding classifier.
All source locations and multiplicities remain binary and the construction
has polynomial source-size and endpoint-bit cost.
\end{proposition}
\begin{proof}
Each local chamber meets at most one region in each source face: the central
face region, triangle arms and quad-chain segments have disjoint chamber
images. An exceptional chamber therefore supplies at most four glued-face
incidences. There are at most $6T$ exceptional chambers, giving the port
bound. Each port has exactly one prism endpoint and one exceptional endpoint.

Clipping removes precisely the exceptional vertices and their incident
arrows from the original face-adjacency graph. All remaining face arrows
are retained through their compressed affine intervals, so their orbits
are exactly the stated prism face-components. A component with no port
has no original adjacency to an exceptional chamber; it is therefore a
whole unmarked cut component. Conversely a whole unmarked component has
no such port. The source record keeps the exact original geometric face
location of every removed core/prism adjacency.

There are only constantly many exceptional points in each incident
tetrahedron and only constantly many affine face pieces per face. Splitting
at their preimages adds only a constant number of pieces per face.
All indices use $O(B+\log(T+1))$ bits; no chamber or repeated face band is
expanded.
\end{proof}

This identifies prism chains attached within a core even when the cut has
only one component. The core cells and their face contacts are local
incidence data, not a completed triangulation of the core. Source edge and
vertex neighborhoods are not supplied by this representation.

\subsection{Weighted port ownership and 0/1-handle profiles}

Give each compressed prism point one cell-count weight. Give each retained
prism/prism face pairing one band-count indicator over its source endpoint
interval. Finally give every interface port a distinct singleton indicator
at its prism point. The dimension is two plus the number of ports.
All work, including overlapping indicators, is charged to the existing
weighted orbit operation.

An orbit's port entries are zero or one. A nonzero port-weight vector can
occur only once, because that particular port belongs to exactly one orbit.
Thus these vectors recover the attached-region port partition without
enumerating points or performing a separate search for every port.
Zero-port rows coalesce repeated unattached regions with binary multiplicities.
For each profile the output records prism-cell count $c$, face-band count $b$
and the ribbon quantity
\[
                       \chi_{01}=c-b.
\]

Here ``ribbon'' refers to the 0/1-handle neighborhood of the prism blocks
and their face bands, with corner attachments kept separate. Its Euler
value follows from its discs and bands. It is not the Euler characteristic
of the full parallelity base: edge handles can fill corner cycles, and
vertex data and boundary patterns still require construction. We do not
identify touching edge strata as a completed bundle, or infer base genus
and boundary counts from $\chi_{01}$.

The number of attached profiles is at most the number of ports. An unattached
profile's cell and band counts are linear functions of its full normal
midsection vector, so their number of distinct values is bounded by the
previous $v+2q$ type bound. This gives a polynomial output with repeated
unattached regions encoded by multiplicity.

\subsection{Independent source clipping and remaining topology}

The API is \code{normal\_prism\_face\_regions}. Its checker independently
constructs affine face segments, walks their exceptional-point events,
reconstructs core contacts and ports, and replays the exact weighted source.
The producer's splitter, summary and orbit planner are not used.
Every source location, port group, cell/band count and ribbon value must
match the independently owned construction. The schema is
\code{normal-prism-interfaces-v1}; large integer fields use the existing
exact hexadecimal transport. Incomplete queries have no structure or proof.
Operation limits mean unverified, and cooperative cancellation propagates.

The source still satisfies the existing finite compact orientable one-torus
contract. The result supplies attached face regions and source-labelled
interfaces inside cores, not their full parallelity topology. Edge/vertex
handles, actual base genus and boundary curves, cyclic ordering of vertical
attachments, core geometry and general later hierarchy-piece validation
remain unresolved. The complete recognizer is not changed and no general
quasi-polynomial recognition bound follows.

\subsection{Exact source incidence and complete-query evidence}
All 1,571 maintained tests pass in 329.583 seconds; the 31 focused tests
pass in 0.908 seconds. The seven new regressions exercise regions inside
primitive cut cores, whole-piece count preservation, empty prism graphs,
self-pairing bands, source/rank/port/partition/weight mutations, shared
verification limits, caps and cancellation. A 16,385-bit source case passes
binary JSON and producer-disabled replay.

The 28.060-second source audit checks 5,289 complete structures: the 5,241
preceding primitive/scaled cases and 48 compatible mixed surfaces. Independent
finite graph expansion agrees with every core-cell record, core/core face
contact, core/prism face location, attached-region port group, cell count,
face-band count and ribbon Euler value. All preceding unattached counts
remain unchanged. Every proof replays with interface, summary and orbit
producers disabled.

The distinction from whole-piece inventory is visible on primitive meridians.
At $T=2,4,8,16$, the prism graph has respectively two, seventeen, 204 and
10,893 points, with two, six, fourteen and thirty source-labelled core-facing
ports. It finds two, three, three and three attached face regions. Each
original cut still has one conservative core and no whole prismatic piece.
Thus the new operator records incidence inside that core, rather than
renaming an earlier whole-piece count.

The audit retains 58 source proofs: 47 source cases with ports, four
primitive inside-core examples and seven large meridian cases. The large
cases cover one-, four- and sixteen-tetrahedron sources at $2^{128}$ or
$2^{500}$ scale, plus a four-tetrahedron $2^{16384}$ case. Source locations,
port ownership and repeated unattached profiles all remain binary. No
expanded graph is constructed for these enormous cases.

Five paired timing rounds measure 80 completed structure calls and sixteen
warmups with A/A controls. Both arms return byte-identical canonical
structures: source core cells and contacts, port locations, region groups
and the complete ribbon profiles. The reference expands the independent
geometric prism graph and its weights. The native arm charges source
validation, geometric clipping, weighted transport, complete independent
replay and result serialization. Neither arm returns missing edge/vertex
handles or full parallelity topology.
\begin{center}
\begin{tabular}{lrrrrr}
\toprule
input & explicit ms & native ms & ratio & old A/A & new A/A\\
\midrule
$T=4$, copies 1 & 0.864 & 3.037 & 0.286 & 0.999 & 0.996\\
$T=16$, copies 1 & 88.150 & 57.157 & 1.299 & 1.047 & 1.018\\
$T=4$, copies 256 & 37.006 & 9.453 & 3.896 & 1.003 & 1.005\\
$T=4$, copies 4096 & 598.350 & 9.252 & 64.801 & 1.018 & 0.998\\
\bottomrule
\end{tabular}
\end{center}


At $T=4$ and 4,096 copies, the paired expansion/native ratio is 64.801,
with median 598.350 versus 9.252 milliseconds. At 256 copies the ratio is
3.896. The small $T=4$ primitive case loses at ratio 0.286. The $T=16$
primitive case has paired ratio 1.299, median 88.150 versus 57.157
milliseconds. That is a source-incidence query gain, not a full recognizer
measurement. Its old A/A ratio is 1.047; timing variability and the small
sample prevent a universal primitive-source speedup claim.

Old A/A ratios range from 0.999 to 1.047 and new ones from 0.996 to 1.018.
All outcomes agree; no capped or unfinished query is counted as a completed
structure. The important encoded-input improvement is avoiding chamber
expansion while retaining exact source interface locations and repetitions.
The complete recognition algorithm and its asymptotic bound remain unchanged.

Runtime release \code{0367cc5c2}, all tests, source incidence audits and
complete-query measurements are retained.

The driver is \code{normal\_orbit\_research.prism\_interfaces}.

Records are under \code{data/prism-interfaces-*}.

Publication review checks 668 current source and evidence hashes, every exact structure/timing outcome and the test
records, and independently replays all 58 saved source proofs with producers
disabled. The next required topology is source edge and vertex handles,
actual parallelity-base boundary curves, and their ordered attachment to the
core. Those are not supplied by the 0/1-handle ribbon profiles.

